\chapter{Mobile Money, Agent Float, and Wallet Unit Economics}
\label{ch:mobile-money-ledgers}
\chapterline

It is 08:30 on Monday morning on a busy street in Nairobi. Grace walks up to a small green corner shop holding a crisp $1{,}000$ Kenyan shilling (\texttt{KES}) banknote, worth roughly eight US dollars. Behind a wire counter sits Joseph, the neighborhood shopkeeper. Joseph sells soap, cooking oil, and sugar, and he also serves as a \textbf{mobile money} agent (\textbf{mobile money} is a phone-based wallet system that stores and sends cash through a mobile \textbf{SIM} -- \emph{Subscriber Identity Module} -- card, with no bank account needed).

Grace hands Joseph her $1{,}000\text{ KES}$ paper bill and reads out her phone number. Joseph taps a short menu code on his rugged agent handset and enters his secret \textbf{PIN} (\emph{Personal Identification Number} passcode). Three seconds later, both phones buzz. Grace's wallet balance just rose by $1{,}000\text{ KES}$. She tucks her phone into her bag and boards the bus to work.

Pause at that wooden counter and ask a physical question. Did Joseph just create $1{,}000\text{ KES}$ of new money out of thin air?

\section{How a Corner-Shop Mobile Money Agent Works}

Joseph cannot print money, and neither can the mobile network. Instead, Joseph runs his kiosk using two separate inventories of his own working capital:
\begin{enumerate}
  \item \textbf{Physical paper cash:} banknotes and coins sitting inside his locked wooden cash drawer.
  \item \textbf{Digital agent float:} pre-purchased electronic balance sitting on his special agent phone (\textbf{agent float} is a shopkeeper's own inventory of digital money held ready to trade with walk-in customers).
\end{enumerate}

Suppose Joseph opens his shop at 08:00 with $20{,}000\text{ KES}$ of paper bills in his drawer and $20{,}000\text{ KES}$ of digital agent float on his phone. His total inventory is $40{,}000\text{ KES}$. Now watch what happens when Grace makes her $1{,}000\text{ KES}$ \textbf{cash-in} (a deposit where a customer trades paper banknotes for digital wallet money):
\begin{itemize}
  \item Joseph places Grace's $1{,}000\text{ KES}$ paper bill into his drawer. His paper cash rises from $20{,}000\text{ KES}$ to $21{,}000\text{ KES}$.
  \item Joseph's agent phone transfers $1{,}000\text{ KES}$ of digital float to Grace's phone. His digital float drops from $20{,}000\text{ KES}$ to $19{,}000\text{ KES}$.
\end{itemize}

Look at Joseph's total inventory after Grace walks away: $21{,}000 + 19{,}000 = 40{,}000\text{ KES}$. Not a single new shilling was created. Joseph simply swapped $1{,}000\text{ KES}$ of his digital money for $1{,}000\text{ KES}$ of Grace's paper money.

Ten minutes later, a carpenter named Peter walks up needing paper cash to buy timber. Peter asks for a $1{,}000\text{ KES}$ \textbf{cash-out} (a withdrawal where a customer trades digital wallet money back into paper banknotes). Peter sends $1{,}000\text{ KES}$ from his phone to Joseph's agent phone, and Joseph hands Peter a $1{,}000\text{ KES}$ paper bill from the drawer. Now Joseph is right back to $20{,}000\text{ KES}$ in paper and $20{,}000\text{ KES}$ in digital float.

\begin{historybox}
In the early 2000s, millions of East Africans owned basic Nokia phones, while fewer than one in five adults had a bank account. Researchers noticed an ingenious everyday habit across Kenya and Uganda. City workers bought scratch cards for prepaid phone airtime and texted the sixteen-digit PIN code to relatives in rural villages. Rural relatives either loaded the calling minutes or sold the unused PIN to a local shopkeeper for paper cash at a steep $15\%$ discount. Prepaid airtime had quietly become makeshift digital cash.

Vodafone and Kenya's telecom operator Safaricom turned that street invention into clean payment plumbing. They built a dedicated one-for-one shilling ledger tied to the SIM card and trained thousands of corner shops to swap paper cash and digital float. Launched in March 2007 as \textbf{M-Pesa} (\emph{pesa} means money in Swahili), the network grew within five years to move sums rivaling half of Kenya's GDP. Soon \textbf{MTN MoMo} in Uganda and Ghana, \textbf{Airtel Money}, and \textbf{Wave} in Senegal and C\^{o}te d'Ivoire scaled the same corner-shop model across Africa.
\end{historybox}

\begin{happeningbox}
Have you ever seen a corner kiosk turn a customer away saying ``no float''? Now you know why. On payday mornings, dozens of workers line up to deposit paper cash. Each deposit adds paper bills to Joseph's drawer and drains digital float from his phone. Once his digital float hits $0\text{ KES}$, Joseph cannot accept another deposit until he walks his stack of paper bills to a bank branch or a wholesale ``super-agent'' to buy fresh digital float. On Friday evenings the problem flips: everyone wants paper cash for the weekend market, draining his wooden drawer while his digital float piles up.
\end{happeningbox}

\section{Savanna Pay and the First Law of Mobile Money}

Now step into the shoes of three software engineers -- Nia, Sam, and Tunde -- who launch a fintech company called \textbf{Savanna Pay}. Their goal is simple. First, build a clean local wallet in Nairobi so families and small merchants can collect and send mobile money cheaply. Next, connect that wallet across borders so senders in London and New York can pay relatives and suppliers across Kenya, Uganda, Ghana, and Nigeria.

On launch week, Savanna Pay raises $5{,}000{,}000\text{ KES}$ (about $38{,}500\text{ USD}$) from early angel investors. That investor money lands in Savanna Pay's corporate checking account to pay salaries, cloud servers, and rent.

Savanna Pay connects its app to Kenya's mobile networks using two everyday telecom tools. Customers with smartphones tap a button that triggers an \textbf{STK push} (a SIM Toolkit pop-up screen that asks the user to type their M-Pesa PIN to approve a payment). Customers with basic button phones dial a \textbf{USSD} code (\textbf{USSD} is a short keypad menu code like \texttt{*384\#} that works over ordinary cellular signals with zero internet data).

At 09:00 on Monday, Grace tries out Savanna Pay. She approves an STK push on her phone to move $1{,}000\text{ KES}$ from her M-Pesa wallet into her new Savanna Pay wallet. Savanna Pay charges $0\text{ KES}$ on deposits so new users join easily. One second later, $1{,}000\text{ KES}$ arrives in Savanna Pay's bank pool, and Grace's Savanna Pay app shows a balance of $1{,}000\text{ KES}$.

Sam looks at the admin dashboard and smiles. ``Our bank balance just rose by $1{,}000\text{ KES}$,'' he says. ``Did we just make our first $1{,}000\text{ KES}$ of revenue?''

Nia shakes her head. ``We made zero revenue. And not one shilling of that $1{,}000\text{ KES}$ belongs to us.''

\subsection{Law 1: Customer Float Is Never Your Money}

When Savanna Pay accepts Grace's $1{,}000\text{ KES}$, it issues \textbf{e-money} (\textbf{electronic money}: digital wallet balances that the issuer promises to redeem one-for-one for national currency at any moment). Unlike a commercial bank, a wallet company is not allowed to lend Grace's deposit out as car loans or mortgages.

Instead, payment regulators enforce a strict legal wall. Every shilling belonging to customers must sit inside a \textbf{safeguarding account}, also called a \textbf{trust account} (a ring-fenced bank account held strictly for customers, legally shielded from the company's own bills and creditors).

Look at Savanna Pay's Balance Sheet right after Grace deposits her $1{,}000\text{ KES}$. Savanna Pay actually holds two completely separate worlds of cash:

\begin{center}
\footnotesize
\setlength{\tabcolsep}{4pt}
\begin{tabular}{lr@{\hspace{1.2em}}lr}
\toprule
\multicolumn{2}{l}{\textbf{Assets (What Savanna Pay Controls)}} & \multicolumn{2}{l}{\textbf{Liabilities \& Equity (Claims)}} \\
\midrule
\emph{Safeguarding Perimeter (Client)} & & \emph{Safeguarding Perimeter (Client)} & \\
Safeguarded Trust Account & $1{,}000\text{ KES}$ & Customer Wallets (Grace) & $1{,}000\text{ KES}$ \\[0.4em]
\emph{Corporate Perimeter (House)} & & \emph{Corporate Perimeter (Owners)} & \\
Corporate Operating Cash & $5{,}000{,}000\text{ KES}$ & Paid-In Share Capital & $5{,}000{,}000\text{ KES}$ \\
 & & Retained Earnings (Profit) & $0\text{ KES}$ \\
\midrule
\textbf{Total Assets} & $\mathbf{5{,}001{,}000\text{ \textbf{KES}}}$ & \textbf{Total Claims} & $\mathbf{5{,}001{,}000\text{ \textbf{KES}}}$ \\
\bottomrule
\end{tabular}
\end{center}

Look at the double-entry journal posting for Grace's deposit (recall from Chapter~1 that debiting an Asset increases it, and crediting a Liability increases it):

\begin{center}
\small
\begin{tabular}{llrr}
\toprule
\textbf{Account} & \textbf{Bucket} & \textbf{Debit (KES)} & \textbf{Credit (KES)} \\
\midrule
Safeguarded Trust Bank Account & Asset ($\uparrow$) & $1{,}000.00$ & \\
Customer Wallet Liability -- Grace & Liability ($\uparrow$) & & $1{,}000.00$ \\
\bottomrule
\end{tabular}
\end{center}

Check what changed. Safeguarded cash went up by $1{,}000\text{ KES}$, and money owed back to Grace went up by $1{,}000\text{ KES}$. Savanna Pay's Equity changed by $0\text{ KES}$. Revenue is $0\text{ KES}$.

\begin{formalbox}
\textbf{The Safeguarding Invariant.} At every instant $t$, let $F_{\text{trust}}(t)$ be the liquid cash sitting inside ring-fenced trust accounts, and let $L_{\text{cust}}(t)$ be the sum of every customer's wallet balance (plus any uncompleted customer transfers). Payment law requires
\[
  F_{\text{trust}}(t) \ge L_{\text{cust}}(t).
\]
Corporate expenses such as cloud servers, salaries, marketing, and office rent may only be paid out of \texttt{Corporate Operating Cash}, never out of $F_{\text{trust}}$.
\end{formalbox}

\begin{trapbox}
\textbf{The Deadly Sin of Commingling.} Early startups sometimes open a single business checking account, drop customer wallet deposits into it, and pay their cloud hosting bill from the same pool. Mixing customer float with company spending money is called \textbf{commingling}. Even if founders promise themselves they will ``put the cash back after next month's fundraising round,'' spending $1\text{ KES}$ of customer float drops $F_{\text{trust}}$ below $L_{\text{cust}}$. Across Kenya, the UK, Europe, and the US, commingling and spending customer float is an immediate license-revoking offense that can carry personal liability for company directors.
\end{trapbox}

\begin{recoverybox}
\textbf{Pause and think:} Suppose Savanna Pay goes viral on Tuesday. Five thousand new users download the app and deposit $2{,}000\text{ KES}$ each, putting $10{,}000{,}000\text{ KES}$ into the trust bank account. Before reading on, answer two questions: how much revenue did Savanna Pay earn from those deposits, and how much of that $10{,}000{,}000\text{ KES}$ can Nia use to pay next month's office rent? Both answers are zero. Every shilling is matched by $10{,}000{,}000\text{ KES}$ of customer wallet liability.
\end{recoverybox}

\subsection{How Savanna Pay Legally Gets Paid: The Daily Fee Sweep}

If Savanna Pay can never touch customer float, how do transaction fees ever reach Savanna Pay's operating bank account? Watch what happens when Grace spends from her wallet.

Suppose Grace first tops up another $100\text{ KES}$ (so her wallet holds $1{,}100\text{ KES}$, backed by $1{,}100\text{ KES}$ in the trust account). At lunch, Grace sends $500\text{ KES}$ to her tailor, Hassan, who also uses Savanna Pay. Savanna Pay charges Grace a $10\text{ KES}$ peer-to-peer transfer fee.

Remember from Chapter~1 that Liabilities and Revenue normally live on the right-hand Credit side. Reducing what we owe Grace is therefore a Debit of $510\text{ KES}$. That Debit is balanced on the Credit side by $+500\text{ KES}$ newly owed to Hassan and $+10\text{ KES}$ of earned fee revenue:

\begin{center}
\small
\begin{tabular}{llrr}
\toprule
\textbf{Account} & \textbf{Bucket} & \textbf{Debit (KES)} & \textbf{Credit (KES)} \\
\midrule
Customer Wallet Liability -- Grace & Liability ($\downarrow$) & $510.00$ & \\
Customer Wallet Liability -- Hassan & Liability ($\uparrow$) & & $500.00$ \\
Transfer Fee Revenue & Revenue / Equity ($\uparrow$) & & $10.00$ \\
\bottomrule
\end{tabular}
\end{center}

Look closely at the trust bank account right after this transfer:
\begin{itemize}
  \item \textbf{Cash inside the Trust Bank Account:} still $1{,}100\text{ KES}$, because no external bank wire moved yet.
  \item \textbf{Total Customer Wallet Liabilities:} Grace now has $1{,}100 - 510 = 590\text{ KES}$, and Hassan has $500\text{ KES}$. Together, customers are owed $590 + 500 = 1{,}090\text{ KES}$.
  \item \textbf{Earned Surplus Inside the Trust Account:} $F_{\text{trust}} - L_{\text{cust}} = 1{,}100 - 1{,}090 = 10\text{ KES}$.
\end{itemize}

Those $10$ shillings are no longer owed to Grace or Hassan. Grace paid them as a fee, extinguishing $10\text{ KES}$ of wallet liability and creating $10\text{ KES}$ of Savanna Pay revenue. Every morning after automated checks verify the balances, Savanna Pay's treasury runs a \textbf{Daily Fee Sweep} (a bank transfer moving only earned net fees from the trust account to the company's operating account). Because both bank accounts are Assets, house cash is debited and trust cash is credited:

\begin{center}
\small
\begin{tabular}{llrr}
\toprule
\textbf{Account} & \textbf{Bucket} & \textbf{Debit (KES)} & \textbf{Credit (KES)} \\
\midrule
Corporate Operating Cash & Asset -- House ($\uparrow$) & $10.00$ & \\
Safeguarded Trust Bank Account & Asset -- Trust ($\downarrow$) & & $10.00$ \\
\bottomrule
\end{tabular}
\end{center}

Now the Trust Account holds $1{,}090\text{ KES}$ against $1{,}090\text{ KES}$ of customer wallets, and Corporate Operating Cash holds $5{,}000{,}010\text{ KES}$, ready to pay cloud bills and salaries legally (Figure~\ref{fig:safeguarding-firewall}).

\begin{figure}[htbp]
\centering
\begin{tikzpicture}[>=Latex, font=\small]
  % Safeguarding side (left)
  \node[font=\bfseries\small, text=sourceblue, align=center] (shead)
    {Safeguarding Perimeter\\[-1pt]{\footnotesize (1:1 Customer Trust)}};

  \node[draw=sourceblue!85!black, thick, fill=white, rounded corners=5pt,
        inner xsep=10pt, inner ysep=8pt, minimum width=4.15cm, align=center,
        below=0.35cm of shead] (sasset)
    {\textbf{Asset: Safeguarded Float}\\[2pt]
     {\footnotesize\textcolor{gray!85!black}{Trust Bank Accounts \&}}\\[-1pt]
     {\footnotesize\textcolor{gray!85!black}{Mobile Money Float Pools}}};

  \node[draw=sourceblue!85!black, thick, fill=white, rounded corners=5pt,
        inner xsep=10pt, inner ysep=8pt, minimum width=4.15cm, align=center,
        below=1.10cm of sasset] (sliab)
    {\textbf{Liability: Customer Wallets}\\[2pt]
     {\footnotesize\textcolor{gray!85!black}{Available Balances +}}\\[-1pt]
     {\footnotesize\textcolor{gray!85!black}{Reserved In-Flight Sends}}};

  \draw[<->, very thick, sourceblue] (sasset.south) --
    node[right=5pt, font=\footnotesize\bfseries, text=sourceblue] (sinv) {$F_{\text{trust}} \ge L_{\text{cust}}$}
    (sliab.north);

  % Corporate side (right)
  \node[draw=targetgreen!85!black, thick, fill=white, rounded corners=5pt,
        inner xsep=10pt, inner ysep=8pt, minimum width=4.15cm, align=center,
        right=3.75cm of sasset] (casset)
    {\textbf{Asset: Operating Cash}\\[2pt]
     {\footnotesize\textcolor{gray!85!black}{Pays Cloud, Staff, Rent,}}\\[-1pt]
     {\footnotesize\textcolor{gray!85!black}{SMS, KYC, \& Taxes}}};

  \node[font=\bfseries\small, text=targetgreen!80!black, align=center,
        above=0.35cm of casset] (chead)
    {Corporate Perimeter\\[-1pt]{\footnotesize (House Cash)}};

  \node[draw=targetgreen!85!black, thick, fill=white, rounded corners=5pt,
        inner xsep=10pt, inner ysep=8pt, minimum width=4.15cm, align=center,
        below=1.10cm of casset] (cequity)
    {\textbf{Equity \& Net Fee Profit}\\[2pt]
     {\footnotesize\textcolor{gray!85!black}{Investor Capital + Earned}}\\[-1pt]
     {\footnotesize\textcolor{gray!85!black}{Fees Minus Rail Costs}}};

  \draw[<->, thick, targetgreen!80!black] (casset.south) --
    node[right=5pt, font=\footnotesize\bfseries, text=targetgreen!80!black] (chouse) {House Books}
    (cequity.north);

  \begin{scope}[on background layer]
    \node[draw=sourceblue, thick, fill=sourceblue!5, rounded corners=6pt,
          inner sep=11pt, fit=(shead)(sasset)(sliab)(sinv)] (sbox) {};
    \node[draw=targetgreen!85!black, thick, fill=targetgreen!5, rounded corners=6pt,
          inner sep=11pt, fit=(chead)(casset)(cequity)(chouse)] (cbox) {};
  \end{scope}

  % Central Legal Firewall
  \coordinate (fwtop) at ($(sbox.north east)!0.5!(cbox.north west) + (0,0.22)$);
  \coordinate (fwbot) at ($(sbox.south east)!0.5!(cbox.south west) - (0,0.22)$);
  \draw[ultra thick, dashed, trapred] (fwbot) -- (fwtop);
  \node[above=2pt of fwtop, font=\footnotesize\bfseries, text=trapred] {Legal Firewall};

  % Daily Earned Fee Sweep arrow across the central firewall gap
  \draw[->, very thick, quotientorange!90!black] (sasset.east) --
    node[above=4pt, font=\footnotesize\bfseries, text=quotientorange!85!black, fill=white, inner xsep=3pt, inner ysep=1.5pt] {Daily Earned}
    node[below=4pt, font=\footnotesize\bfseries, text=quotientorange!85!black, fill=white, inner xsep=3pt, inner ysep=1.5pt] {Fee Sweep}
    (casset.west);
\end{tikzpicture}
\caption{The legal wall between customer float and house cash. Customer deposits stay locked inside the safeguarding perimeter; only reconciled net fees cross the wall via the daily fee sweep.}
\label{fig:safeguarding-firewall}
\end{figure}

\section{Unit Economics: Where Every Shilling of Fee Goes}

Why was Grace's morning cash-in free, while sending or withdrawing money carries a fee? Mobile money networks price the lifecycle of a shilling carefully:
\begin{itemize}
  \item \textbf{Cash-in is usually free to the customer:} You want paper cash to enter the digital network with zero friction. Even so, cash-in is not free to the operator: if Grace deposits at Joseph's kiosk, the network pays Joseph a small commission for taking her paper bill, or Savanna Pay absorbs a tiny STK-push gateway fee.
  \item \textbf{Peer-to-peer (P2P) wallet transfers carry a tiny fee:} Moving balances inside Savanna Pay's own database costs almost nothing in external bank fees, so a $0.5\%$ to $1.0\%$ fee is high-margin revenue.
  \item \textbf{Cash-out carries a tiered fee split across partners:} Pulling digital money back out into paper cash or an external mobile network is expensive because corner-shop agents and telecom rails must be paid.
\end{itemize}

Figure~\ref{ch07:fig-wallet-lifecycle} traces \$100 through cash-in, peer-to-peer transfer, and cash-out to show why trust float covers wallet liabilities at every step.

\begin{figure}[htbp]
\centering
\begin{tikzpicture}[>=Latex, font=\footnotesize]
  \tikzset{
    actbox/.style={draw=gray!80!black, thick, fill=gray!6, rounded corners=5pt,
      inner xsep=9pt, inner ysep=7pt, minimum width=2.85cm, align=left},
    floatbox/.style={draw=sourceblue!85!black, thick, fill=sourceblue!6, rounded corners=5pt,
      inner xsep=9pt, inner ysep=7pt, minimum width=3.35cm, align=left},
    liabbox/.style={draw=targetgreen!85!black, thick, fill=targetgreen!6, rounded corners=5pt,
      inner xsep=9pt, inner ysep=7pt, minimum width=4.15cm, align=left}
  }

  % Row 1: Cash-In
  \node[actbox] (a1) {\textbf{1. Cash-In (\$100)}\\[2pt]Alice deposits \$100 cash};
  \node[floatbox, right=0.60cm of a1] (f1) {\textbf{Trust Float: $+$\$100.00}\\[2pt]Pool holds \textbf{\$100.00}};
  \node[liabbox, right=0.75cm of f1] (l1) {\textbf{Alice Wallet: $+$\$100.00}\\[2pt]Customer liability $= \mathbf{\$100.00}$};

  % Row 2: P2P Send
  \node[actbox, below=0.65cm of a1] (a2) {\textbf{2. P2P Send (\$40)}\\[2pt]Alice pays Bob (\$0.50 fee)};
  \node[floatbox, right=0.60cm of a2] (f2) {\textbf{Trust Float: \$0.00}\\[2pt]Pool stays at \textbf{\$100.00}};
  \node[liabbox, right=0.75cm of f2] (l2) {\textbf{Alice $-$\$40 $\to$ Bob $+$\$39.50}\\[2pt]Wallets \textbf{\$99.50} $+$ Fee \textbf{\$0.50}};

  % Row 3: Cash-Out
  \node[actbox, below=0.65cm of a2] (a3) {\textbf{3. Cash-Out (\$39.50)}\\[2pt]Bob withdraws via agent};
  \node[floatbox, right=0.60cm of a3] (f3) {\textbf{Trust Float: $-$\$39.50}\\[2pt]Pool holds \textbf{\$60.50}};
  \node[liabbox, right=0.75cm of f3] (l3) {\textbf{Bob: $-$\$39.50 (to \$0)}\\[2pt]Alice \textbf{\$60.00} $+$ Fee \textbf{\$0.50}};

  % Column headers
  \node[above=0.25cm of a1, font=\footnotesize\bfseries, text=gray!80!black] {User Stage};
  \node[above=0.25cm of f1, font=\footnotesize\bfseries, text=sourceblue!90!black] {Trust Float Asset ($F_{\text{trust}}$)};
  \node[above=0.25cm of l1, font=\footnotesize\bfseries, text=targetgreen!85!black] {Wallets ($L_{\text{cust}}$) + Earned Fee};

  % Arrows
  \foreach \r in {1,2,3} {
    \draw[->, thick, gray!75!black] (a\r.east) -- (f\r.west);
    \draw[<->, thick, sourceblue!85!black] (f\r.east) -- node[above=2pt, font=\scriptsize\bfseries, text=sourceblue!90!black] {$\ge$} (l\r.west);
  }
  \draw[->, thick, gray!65!black] (a1.south) -- (a2.north);
  \draw[->, thick, gray!65!black] (a2.south) -- (a3.north);

  % Bottom invariant bar centered beneath the 3 columns
  \coordinate (midrow3) at ($(a3.south west)!0.5!(l3.south east)$);
  \node[draw=structurepurple!85!black, thick, fill=structurepurple!7, rounded corners=5pt,
        inner xsep=11pt, inner ysep=6pt, font=\footnotesize\bfseries, text=structurepurple!90!black,
        anchor=north] at ([yshift=-0.38cm]midrow3)
    {Invariant at Every Millisecond: $F_{\text{trust}}(t) \ge L_{\text{cust}}(t)$ (Earned \$0.50 sweeps to House Cash)};
\end{tikzpicture}
\caption{Three-stage mobile money balance conservation across cash-in, peer-to-peer send, and cash-out. Trust float rises and falls in lockstep with customer liabilities.}
\label{ch07:fig-wallet-lifecycle}
\end{figure}

Let us trace one external withdrawal to the exact shilling. Suppose Hassan cashes out $1{,}000\text{ KES}$ from Savanna Pay through a partner agent rail. Savanna Pay's tiered fee table charges Hassan a $28\text{ KES}$ withdrawal fee, reducing his wallet by $1{,}028\text{ KES}$ (assuming Hassan held at least $1{,}028\text{ KES}$). Where does that $28\text{ KES}$ go?
\begin{enumerate}
  \item \textbf{Corner-shop agent commission ($10\text{ KES}$):} Paid to the retail shopkeeper so they keep stocked drawers of paper banknotes on the street.
  \item \textbf{Telecom payout rail tariff ($6\text{ KES}$):} Deducted by the mobile network switch for delivering the payout.
  \item \textbf{One-time password (OTP) texts and KYC verification ($2\text{ KES}$):} Paid by Savanna Pay for \textbf{OTP} (\emph{one-time password}) login codes sent via \textbf{SMS} (\emph{Short Message Service} text messages) and amortized \textbf{KYC} (\emph{Know Your Customer}) identity verification of government IDs and selfies.
  \item \textbf{Savanna Pay net contribution ($10\text{ KES}$):} Out of the $28\text{ KES}$ fee Hassan saw on screen, Savanna Pay keeps $28 - 10 - 6 - 2 = 10\text{ KES}$.
\end{enumerate}

Look at the balanced double-entry posting when Hassan's $1{,}000\text{ KES}$ cash-out settles (where the external agent/telecom pool deducts the $1{,}000\text{ KES}$ payout plus $10\text{ KES}$ agent commission plus $6\text{ KES}$ rail tariff $= 1{,}016\text{ KES}$, and $2\text{ KES}$ accrues to Savanna Pay's SMS vendor):

\begin{center}
\small
\begin{tabular}{llrr}
\toprule
\textbf{Account} & \textbf{Bucket} & \textbf{Debit (KES)} & \textbf{Credit (KES)} \\
\midrule
Customer Wallet Liability -- Hassan & Liability ($\downarrow$) & $1{,}028.00$ & \\
Agent Commission Expense & Expense ($\uparrow$) & $10.00$ & \\
Telecom Rail Tariff Expense & Expense ($\uparrow$) & $6.00$ & \\
SMS \& KYC Verification Expense & Expense ($\uparrow$) & $2.00$ & \\
Safeguarded Float -- Agent/Telco Pool & Asset ($\downarrow$) & & $1{,}016.00$ \\
Withdrawal Fee Revenue & Revenue ($\uparrow$) & & $28.00$ \\
Accounts Payable -- SMS Vendor & Liability -- House ($\uparrow$) & & $2.00$ \\
\bottomrule
\end{tabular}
\end{center}

Every side sums to $1{,}046.00\text{ KES}$. Inside the safeguarded float pool, customer liability fell by $1{,}028\text{ KES}$ while cash fell by $1{,}016\text{ KES}$, leaving $12\text{ KES}$ of corporate surplus to sweep across the wall (from which Savanna Pay later pays the $2\text{ KES}$ SMS bill, netting $10\text{ KES}$).

\subsection{Basis Points, GTV, and the Twenty-Million-Dollar Illusion}

Because payment fees are small slices of $1\%$, payment teams measure every fee and cost in \textbf{basis points}. One \textbf{basis point} (written \textbf{1 bp}) is one hundredth of one percent, or one cent on every hundred dollars:
\[
  1\text{ bp} = 0.01\% = 0.0001, \qquad 100\text{ bps} = 1.00\%, \qquad 10{,}000\text{ bps} = 100\%.
\]
Using basis points prevents costly decimal-point typos. Subtracting $110\text{ bps}$ of rail costs from $145\text{ bps}$ of fees immediately gives $35\text{ bps}$, with no strings of leading zeros.

Payment companies also separate \textbf{Gross Transaction Volume (GTV)} -- the total customer money moved across the platform in a period, also called Total Payment Volume -- from actual company revenue. Dividing fee revenue (or fee revenue minus direct rail costs) by GTV gives the fraction kept per dollar moved, and multiplying by $10{,}000$ expresses that fraction in basis points:
\begin{align*}
  \text{Gross Take Rate (bps)} &= 10{,}000 \times \frac{\text{Gross Fee Revenue}}{\text{GTV}}, \\
  \text{Net Margin (bps)} &= 10{,}000 \times \frac{\text{Gross Fee Revenue} - \text{Direct Rail Costs}}{\text{GTV}}.
\end{align*}

\begin{intuitionbox}
Think of Savanna Pay as a freight railway. \textbf{GTV} is the market value of the cargo inside customers' shipping containers. \textbf{Gross Fee Revenue} is the freight ticket customers pay to move the containers. \textbf{Direct Rail Costs} are diesel fuel and track tolls. No railway accountant ever records the value of customers' televisions and coffee beans as railway revenue. In payments, the cargo happens to be money itself, which tricks inexperienced founders into confusing $20{,}000{,}000\text{ USD}$ of cargo with $20{,}000{,}000\text{ USD}$ of revenue.
\end{intuitionbox}

Fast-forward eighteen months. Savanna Pay is growing fast, moving $\mathbf{20{,}000{,}000\text{ \textbf{USD}}}$ per month in GTV across $400{,}000$ transfers (an average transfer of $50.00\text{ USD}$). On an investor slide, ``$240\text{M USD}$ annualized payment volume'' sounds huge. Table~\ref{tab:payment-unit-econ} shows what actually reaches Savanna Pay's operating bank account each month after paying gateway fees, agent commissions, one-time password (OTP) texts, Know Your Customer (KYC) identity verification, and fraud losses.

\begin{table}[htbp]
\centering
\caption{Savanna Pay monthly unit-economics waterfall on $20{,}000{,}000\text{ USD}$ GTV across $400{,}000$ transfers ($50.00\text{ USD}$ average transfer).}
\label{tab:payment-unit-econ}
\footnotesize
\setlength{\tabcolsep}{4.0pt}
\begin{tabular}{lrrr}
\toprule
\textbf{Line Item} & \textbf{Rate (bps)} & \textbf{Per \$50 Send (\$)} & \textbf{Monthly Total (\$)} \\
\midrule
\textbf{Gross Transaction Volume (GTV)} & $10{,}000\text{ bps}$ & \$50.000 & \$20,000,000 \\
\midrule
Customer fees \& FX spread (\textbf{Gross Revenue}) & $145\text{ bps}$ & \$0.725 & \$290,000 \\
\quad Less: Pay-in gateway / STK push / card fees & $-48\text{ bps}$ & $-$\$0.240 & $-$\$96,000 \\
\quad Less: Agent commissions \& payout rail tariffs & $-37\text{ bps}$ & $-$\$0.185 & $-$\$74,000 \\
\quad Less: OTP SMS, KYC checks, \& compliance APIs & $-13\text{ bps}$ & $-$\$0.065 & $-$\$26,000 \\
\quad Less: Fraud, reversals, \& unrecovered disputes & $-12\text{ bps}$ & $-$\$0.060 & $-$\$24,000 \\
\midrule
\textbf{Total Direct Variable Costs (Payment COGS)} & $-110\text{ bps}$ & $-$\$0.550 & $-$\$220,000 \\
\midrule
\textbf{Net Contribution Margin} & $\mathbf{35}\textbf{ bps}$ & \textbf{\$0.175} & \textbf{\$70,000} \\
\bottomrule
\end{tabular}
\end{table}

Read the bottom row of Table~\ref{tab:payment-unit-econ} carefully. Out of $20{,}000{,}000\text{ USD}$ moving through the pipes, Savanna Pay keeps \textbf{17.5 cents} per $50\text{ USD}$ transfer, or $\mathbf{70{,}000\text{ \textbf{USD}}}$ \textbf{per month}. That $70{,}000\text{ USD}$ must pay every engineer, compliance officer, server database, audit fee, and office lease. If fixed monthly salaries and servers cost $95{,}000\text{ USD}$, Savanna Pay is still burning $25{,}000\text{ USD}$ of investor cash every month.

\begin{walkbox}
\textbf{Test your intuition:} Suppose Savanna Pay's marketing lead proposes a ``Half-Price Month'' promotion that cuts customer fees by $45\text{ bps}$ (from $145\text{ bps}$ down to $100\text{ bps}$), promising that monthly volume will double from $20\text{M USD}$ to $40\text{M USD}$. Should Nia approve it? Look at the variable costs: each transfer still costs $110\text{ bps}$ in rails, agents, and SMS. At $100\text{ bps}$ of revenue, net margin turns negative ($-10\text{ bps}$). Doubling volume to $40\text{M USD}$ would lose $40{,}000\text{ USD}$ on transfers alone before paying a single salary.
\end{walkbox}

\section{The Classic Engineering Bug: Race Conditions on Payouts}

Knowing the accounting rules on paper is half the job. Your backend database has to enforce those exact rules when mobile networks drop packets and users tap buttons twice.

During Savanna Pay's beta test, Tunde notices a scary anomaly in the staging logs. A tester named Kofi held $1{,}050\text{ KES}$ in his wallet. Sitting on a bumpy bus with patchy 3G signal, Kofi tapped ``Withdraw $1{,}000\text{ KES}$ to M-Pesa'' (with a $28\text{ KES}$ fee, totaling $1{,}028\text{ KES}$). Nothing appeared to happen for half a second, so Kofi tapped the button again.

Look at how a naive backend service handles those two taps when it stores a single \texttt{users.balance} column and calls the external M-Pesa API \emph{before} updating the database:
\begin{enumerate}
  \item \textbf{Millisecond 0:} Request~A queries Kofi's row (\texttt{WHERE id = 42}). It reads $1{,}050\text{ KES} \ge 1{,}028\text{ KES}$, so the check passes.
  \item \textbf{Millisecond 9:} Request~B arrives before M-Pesa responds to Request~A. Request~B runs the exact same \texttt{SELECT}, also sees $1{,}050\text{ KES} \ge 1{,}028\text{ KES}$, and also passes.
  \item \textbf{Millisecond 400:} Both threads call Safaricom's external M-Pesa payout API. Both external payouts succeed. M-Pesa sends $2{,}000\text{ KES}$ of real cash out of Savanna Pay's float pool to Kofi's phone, even though Kofi only owned $1{,}050\text{ KES}$.
\end{enumerate}

Savanna Pay just lost $978\text{ KES}$ out of its safeguarded pool, knocking $F_{\text{trust}}$ below $L_{\text{cust}}$ in under half a second.

\begin{trapbox}
\textbf{The Mirror Trap: Auto-Refunding on an HTTP Timeout.} Suppose Sam tries a quick fix: deduct $1{,}028\text{ KES}$ from \texttt{users.balance} \emph{first}, call the M-Pesa API, and if the HTTP call throws an error, add $1{,}028\text{ KES}$ right back. That creates an equally dangerous exploit. Suppose M-Pesa receives Request~A and pays Kofi $1{,}000\text{ KES}$, while a proxy drop causes Savanna Pay's HTTP client to throw \texttt{504 Gateway Timeout} after $30$ seconds. If your \texttt{catch} block automatically refunds $1{,}028\text{ KES}$, Kofi keeps the first $1{,}000\text{ KES}$ on M-Pesa \emph{and} withdraws the refunded $1{,}028\text{ KES}$ a second time. In distributed payments, a network timeout means \textbf{outcome unknown}, never \textbf{payment failed}.
\end{trapbox}

\subsection{The Safe Pattern: Reserve, External Payout, then Finalize}

Production mobile-money switches prevent both bugs using three simple engineering rules:
\begin{enumerate}
  \item \textbf{Store integer minor units only:} Never store money in binary \texttt{FLOAT} or \texttt{DOUBLE} (where $0.10 + 0.20 \neq 0.30$). Store every balance as a 64-bit integer (\texttt{BIGINT}) counting minor units such as cents ($\text{KES } 1{,}028.00 = \texttt{102800}$).
  \item \textbf{Attach an idempotency key to every step.} An \textbf{idempotency key} is a unique tag created per user action. If the phone or server retries the same request five times, the database executes it once and returns the saved receipt on repeats.
  \item \textbf{Split each wallet into \texttt{Available} and \texttt{In-Flight Reserved}:} Model every customer's liability as two sub-accounts so funds lock locally \emph{before} any network call leaves your server.
\end{enumerate}

Trace Kofi's $1{,}000\text{ KES}$ withdrawal ($28\text{ KES}$ fee, $16\text{ KES}$ external agent/telco tariff) through this three-step state machine (Figure~\ref{fig:two-phase-wallet}).

\paragraph{Step 1: Reserve Locally Before Calling the Network}
Inside one atomic SQL transaction that locks Kofi's row (\texttt{SELECT ... FOR UPDATE}), verify \texttt{Available} $\ge 102800$. Move $1{,}028.00\text{ KES}$ from \texttt{Customer Available} to \texttt{Customer In-Flight Reserved} and commit \emph{before} calling M-Pesa:

\begin{center}
\small
\begin{tabular}{llrr}
\toprule
\textbf{Step 1 (\texttt{RESERVE})} & \textbf{Bucket} & \textbf{Debit (KES)} & \textbf{Credit (KES)} \\
\midrule
Customer Available -- Kofi & Liability ($\downarrow$) & $1{,}028.00$ & \\
Customer In-Flight Reserved -- Kofi & Liability ($\uparrow$) & & $1{,}028.00$ \\
\bottomrule
\end{tabular}
\end{center}

When Kofi's second tap arrives $9\text{ ms}$ later, either it carries the same idempotency key (returning the in-progress status harmlessly) or it sees \texttt{Available} $= 22.00\text{ KES} < 1{,}028.00\text{ KES}$ and rejects cleanly before touching M-Pesa.

\paragraph{Step 2a: Finalize on Confirmed External Receipt}
Once M-Pesa confirms delivery with a receipt ID, post the \texttt{SETTLE} batch (keyed uniquely by \texttt{pay\_id:SETTLE} so duplicate webhooks are ignored):

\begin{center}
\small
\begin{tabular}{llrr}
\toprule
\textbf{Step 2a (\texttt{SETTLE})} & \textbf{Bucket} & \textbf{Debit (KES)} & \textbf{Credit (KES)} \\
\midrule
Customer In-Flight Reserved -- Kofi & Liability ($\downarrow$) & $1{,}028.00$ & \\
Agent \& Telco Rail Expense & Expense ($\uparrow$) & $16.00$ & \\
Safeguarded Float -- M-Pesa Pool & Asset ($\downarrow$) & & $1{,}016.00$ \\
Withdrawal Fee Revenue & Revenue ($\uparrow$) & & $28.00$ \\
\bottomrule
\end{tabular}
\end{center}

\paragraph{Step 2b: Hold on Timeout, Reverse Only on Confirmed Failure}
If the HTTP call times out, move the payment status to \texttt{PENDING\_RECON} while keeping Kofi's $1{,}028.00\text{ KES}$ locked inside \texttt{In-Flight Reserved}. A background worker queries M-Pesa's transaction-status API using your original idempotency key. Only if M-Pesa definitively confirms that the payout failed (for instance, an unregistered phone number) does your worker post \texttt{REVERSE}, moving $1{,}028.00\text{ KES}$ back from \texttt{In-Flight Reserved} to \texttt{Customer Available}.

\begin{figure}[htbp]
\centering
\begin{tikzpicture}[>=Latex, font=\footnotesize]
  \tikzset{
    state/.style={draw=sourceblue!85!black, thick, fill=sourceblue!7, rounded corners=5pt,
      inner xsep=9pt, inner ysep=7pt, align=center},
    termok/.style={draw=targetgreen!85!black, thick, fill=targetgreen!10, rounded corners=5pt,
      inner xsep=9pt, inner ysep=7pt, align=center},
    termrev/.style={draw=quotientorange!90!black, thick, fill=quotientorange!10, rounded corners=5pt,
      inner xsep=9pt, inner ysep=7pt, align=center},
    ambig/.style={draw=trapred!85!black, thick, fill=trapred!8, rounded corners=5pt,
      inner xsep=9pt, inner ysep=7pt, align=center}
  }

  \node[state] (init) {\textbf{INITIATED}\\[1pt]{\scriptsize App Tap + Key}};
  \node[state, right=2.50cm of init] (res) {\textbf{RESERVED}\\[1pt]{\scriptsize Dr Avail / Cr In-Flight}};
  \node[ambig, right=2.35cm of res] (pend) {\textbf{PENDING\_RECON}\\[1pt]{\scriptsize Timeout: Keep Lock}};
  \node[termok, above=1.40cm of pend] (set) {\textbf{SETTLED}\\[1pt]{\scriptsize Dr In-Flight / Cr Float}};
  \node[termrev, below=1.40cm of pend] (rev) {\textbf{REVERSED}\\[1pt]{\scriptsize Dr In-Flight / Cr Avail}};

  % Orthogonal transitions with zero crossings
  \draw[->, thick, sourceblue!90!black] (init.east) --
    node[above=3pt, font=\scriptsize\bfseries] {Row Lock OK} (res.west);

  \draw[->, thick, trapred!85!black] (res.east) --
    node[above=3pt, font=\scriptsize\bfseries, text=trapred!85!black] {504 Timeout} (pend.west);

  \draw[->, thick, targetgreen!80!black, rounded corners=5pt] (res.north) |- (set.west)
    node[pos=0.74, above=3pt, font=\scriptsize\bfseries, text=targetgreen!80!black] {Receipt Webhook};

  \draw[->, thick, quotientorange!85!black, rounded corners=5pt] (res.south) |- (rev.west)
    node[pos=0.74, below=3pt, font=\scriptsize\bfseries, text=quotientorange!85!black] {Confirmed Reject};

  \draw[->, thick, dashed, targetgreen!80!black] (pend.north) --
    node[left=5pt, font=\scriptsize\bfseries, text=targetgreen!80!black] {Status: Paid} (set.south);

  \draw[->, thick, dashed, quotientorange!85!black] (pend.south) --
    node[left=5pt, font=\scriptsize\bfseries, text=quotientorange!85!black] {Status: Failed} (rev.north);
\end{tikzpicture}
\caption{Safe three-step payout state machine. Funds lock in \texttt{RESERVED} before any external API call, and timeouts hold inside \texttt{PENDING\_RECON} until the external rail confirms whether cash left the pool.}
\label{fig:two-phase-wallet}
\end{figure}

Three compact SQL tables implement this entire double-entry engine:

\begin{lstlisting}[language=SQL]
CREATE TABLE ledger_accounts (
  id             BIGSERIAL PRIMARY KEY,
  code           TEXT NOT NULL UNIQUE,      -- e.g. 'WALLET_AVAIL:user_42'
  currency       CHAR(3) NOT NULL,          -- 'KES', 'UGX', 'GHS', 'NGN', 'USD'
  balance_minor  BIGINT NOT NULL DEFAULT 0  -- integer cents; locked FOR UPDATE
);

CREATE TABLE journal_batches (
  id               BIGSERIAL PRIMARY KEY,
  idempotency_key  TEXT NOT NULL UNIQUE,    -- e.g. 'pay_981:RESERVE'
  created_at       TIMESTAMPTZ NOT NULL DEFAULT now(),
  memo             TEXT NOT NULL
);

CREATE TABLE ledger_postings (
  id             BIGSERIAL PRIMARY KEY,
  batch_id       BIGINT NOT NULL REFERENCES journal_batches(id),
  account_id     BIGINT NOT NULL REFERENCES ledger_accounts(id),
  currency       CHAR(3) NOT NULL,
  signed_minor   BIGINT NOT NULL CHECK (signed_minor <> 0)
  -- Convention: Debit = +signed_minor, Credit = -signed_minor.
  -- Trigger checks SUM(signed_minor) = 0 for every (batch_id, currency).
);
\end{lstlisting}

Because Debits are positive integers and Credits are negative integers, every journal batch sums to zero (\texttt{SUM(signed\_minor) = 0}). Inside one country and one currency, Savanna Pay's wallet is now rock solid. In Chapter~\ref{ch:fx-corridors}, we ask what happens when a sender in London taps Send in British pounds and Grace's phone in Nairobi buzzes ten seconds later with Kenyan shillings.

\section*{Exercises}
\addcontentsline{toc}{section}{Exercises}

\begin{exercise}[Agent Float vs Paper Drawer Inventory]
Joseph opens his Nairobi kiosk at 08:00 with $15{,}000\text{ KES}$ in paper banknotes and $25{,}000\text{ KES}$ in digital agent float on his handset.
\begin{enumerate}
  \item Between 08:00 and 11:00, twelve commuters deposit $2{,}000\text{ KES}$ each (cash-in), and three shoppers withdraw $1{,}000\text{ KES}$ each (cash-out). Compute Joseph's paper drawer balance and digital agent float at 11:00 (ignoring commissions).
  \item At 11:05, a merchant arrives to deposit $6{,}000\text{ KES}$ in paper cash. Can Joseph serve her immediately? Explain physically what happens and what Joseph must do at the bank branch to rebalance his kiosk.
\end{enumerate}
\end{exercise}

\begin{exercise}[Safeguarding Wall and the Daily Fee Sweep]
Savanna Pay starts Monday with $3{,}000{,}000\text{ KES}$ of investor capital in \texttt{Corporate Operating Cash} and $0\text{ KES}$ in its \texttt{Safeguarded Trust Account}. During Monday:
\begin{enumerate}
  \item $500$ users deposit $4{,}000\text{ KES}$ each via free bank and STK transfers into the trust account.
  \item Users send $400$ wallet-to-wallet transfers of $1{,}000\text{ KES}$ each, paying a $10\text{ KES}$ fee per transfer.
  \item Users make $200$ external M-Pesa withdrawals of $2{,}000\text{ KES}$ each, paying a $40\text{ KES}$ fee per withdrawal, while the external rail deducts an $18\text{ KES}$ tariff per withdrawal directly from the trust pool.
  \item Savanna Pay pays a $90{,}000\text{ KES}$ cloud server invoice.
\end{enumerate}
First, compute total \texttt{Customer Wallet Liability} and actual trust cash at midnight before any sweep. Second, compute the exact legal \textbf{Daily Fee Sweep} and ending \texttt{Corporate Operating Cash} after paying the cloud bill and running the sweep.
\end{exercise}

\begin{exercise}[Unit Economics in Basis Points]
Savanna Pay's Uganda wallet processes $12{,}000{,}000$ \texttt{USD} in monthly GTV with an average transfer size of $30.00$ \texttt{USD}. Customers pay an average gross take rate of $135\text{ bps}$. Variable costs per transfer are: pay-in gateway fee $45\text{ bps}$, outbound payout tariff $0.09$ \texttt{USD} flat per transfer, SMS and KYC checks $0.03$ \texttt{USD} flat per transfer, and fraud losses $8\text{ bps}$. Fixed monthly salaries and servers cost $55{,}000$ \texttt{USD}.
\begin{enumerate}
  \item Convert every flat dollar cost into basis points at the $30.00$ \texttt{USD} average ticket, and calculate monthly Net Contribution Margin in both basis points and dollars.
  \item Is the Uganda wallet cash-flow positive after fixed costs? Compute the exact monthly GTV needed to break even at the same ticket size and fee schedule.
\end{enumerate}
\end{exercise}

\begin{exercise}[Race Condition vs Two-Phase Reservation]
Grace holds \texttt{350000} shilling cents ($3{,}500.00$ \texttt{KES}) in her Savanna Pay wallet. Two concurrent requests arrive $6\text{ ms}$ apart with different client keys, each asking to withdraw $2{,}500.00$ \texttt{KES} plus a $35.00$ \texttt{KES} fee (\texttt{253500} minor units) to M-Pesa.
\begin{enumerate}
  \item Trace the exact database reads, external API calls, and final cash loss under a buggy handler that checks \texttt{balance} $\ge 253500$ without row locking and deducts \texttt{balance} inside the webhook callback.
  \item Trace the same two requests under the \texttt{Available} $\to$ \texttt{In-Flight Reserved} pattern with row locking (\texttt{FOR UPDATE}), including what happens when Request~1 hits a \texttt{504} timeout and M-Pesa delivers the receipt webhook ten minutes later.
\end{enumerate}
\end{exercise}
